\section{Die Einheit}

\textbf{[K]} Die elementare Informationseinheit dieser Theorie ist das
\emph{Windungsquant}
\[
  \Delta w = 1 .
\]

Es ist der kleinste Unterschied zwischen zwei Windungsklassen auf dem $T^4$.
Ein halbes Windungsquant gibt es nicht -- die Fundamentalgruppe
$\pi_1(T^4)\cong\mathbb{Z}^4$ ist ganzzahlig (A030).

\textbf{[B]} Drei Eigenschaften folgen unmittelbar, und sie sind der Grund,
warum diese Wahl gegenüber jeder zellbasierten Definition vorzuziehen ist:

\begin{itemize}
  \item \textbf{ganzzahlig} -- keine willkürliche Feinheit;
  \item \textbf{koordinatenfrei} -- unabhängig von Ursprung, Orientierung und
    Parametrisierung;
  \item \textbf{energieunabhängig} -- die Einheit ist dieselbe auf jeder Skala,
    weil sie topologisch und nicht metrisch ist.
\end{itemize}

Die letzte Eigenschaft ist die stärkste. Eine Informationseinheit, die von der
Energie abhinge, wäre keine Einheit, sondern eine Skalenwahl.

\section{Warum keine Elementarzellen}

\textbf{[B]} Die naheliegende Alternative -- den Raum in Zellen zu teilen und
Information zu zählen -- scheitert an denselben Einwänden, die in A030 aufgezählt
sind: Ursprung, Orientierung, Koordinatenabhängigkeit, und eine Entropie, die
von einer willkürlichen Einteilung abhängt.

Beim Windungsquant laufen alle diese Einwände leer, weil eine Windungsklasse
keinen Ort hat. Sie ist die Eigenschaft eines geschlossenen Weges, nicht eines
Gebiets. Und der Torus ist randlos, $\partial T^4=\emptyset$ -- eine
flächenbasierte Definition scheidet ebenfalls aus, weil keine Fläche
ausgezeichnet ist.

\section{Betrag und Phase}

\textbf{[K]} Der Massenoperator zerfällt in zwei Kanäle (A070, A120): Beträge
und Phasen der Spektralkoeffizienten. Für den Informationsbegriff ist die
Trennung wesentlich.

\textbf{[B]} Ein Betragsspektrum legt die Phase nur bis auf einen
Allpass-Faktor fest. Es gibt daher Information, die im Betragskanal
\emph{prinzipiell} nicht sichtbar ist -- der Nachweis dazu steht in A120, wo
gezeigt wird, dass $\theta$ von keinem symmetrischen, zählenden Funktional
erfasst wird.

\textbf{[SETZUNG]} Daraus folgt eine Buchungsregel für die ganze Serie: eine
Aussage, die nur Beträge verwendet, kann über Phaseninformation nichts
behaupten -- weder ihre Existenz noch ihre Abwesenheit.

\section{Der Bezug zu Landauer}

\textbf{[B]} Das Landauer-Prinzip verknüpft das Löschen eines Bits mit einer
Mindestenergie $k_B T\ln 2$. In der Sprache dieser Theorie ist das eine
Aussage der \emph{Übersetzungsschicht}: sie enthält $k_B$ und $T$, also
zwei Größen, die nach A080 Brücke und Verkleidung sind.

Das Windungsquant liegt darunter. Es ist die topologische Einheit; die
Landauer-Energie ist, was sie in einer bestimmten thermischen Einkleidung
kostet.

\textbf{[B]} \textbf{Die Ableitung.} Auf dem 4-Torus ist
$\pi_1(T^4)\cong\mathbb Z^4$: die erste Homotopiegruppe ist
$\mathbb Z^4$, jede Richtung trägt diskrete Windungszahlen. Das Löschen eines
Bits entspricht dem Übergang $\Delta w = 1$ -- einer Änderung der Windungszahl
um eins in einer Richtung. Dieser Übergang ist irreversibel (Typ-II-Verlust,
A090): die Windungszahl ist eine topologische Invariante, ihr Verlust kann nicht
rückgängig gemacht werden ohne Energieaufwand.

Der Energieaufwand folgt aus der thermodynamischen Einkleidung dieses
topologischen Übergangs. Ein System im thermischen Gleichgewicht bei Temperatur
$T$ muss für einen irreversiblen Bit-Übergang mindestens die Entropie
$\Delta S = k_B \ln 2$ abführen (zwei Zustände, einer gelöscht). Die
Mindestenergie ist:
\[
  E_L = T\cdot\Delta S = k_B T\ln 2.
\]
Das Windungsquant $\Delta w=1$ ist der geometrische Träger dieses Schritts; die
Landauer-Energie ist seine thermodynamische Verkleidung. Darunter: ein
topologischer Übergang. Darüber: ein thermischer Energiebetrag.

\textbf{[K]} \textbf{Grenzen.} Die Ableitung setzt voraus, dass der Bit-Übergang
mit dem Windungsquant identifiziert werden darf -- das ist die physikalische
Behauptung, die nicht weiter reduziert wird. Die Landauer-Energie $k_BT\ln2$
ist eine Untergrenze; tatsächliche Löschvorgänge kosten mehr. Ob genau $\Delta
w=1$ die einzige Möglichkeit ist, einen Bit zu löschen, ist die zweite offene Frage
(s.\ Was offen bleibt).

\section{Der Umkehrtest}

\textbf{[B]} Die Projektionskette (A090) ist in einer Richtung verlustfrei und
in der anderen nicht. Für den Informationsbegriff heißt das: die Frage
\enquote{wie viel Information geht verloren?} hat nur bei Typ~I und Typ~II eine
Antwort, bei Typ~III lautet sie null.

Beim Ausrollen (Typ~I) ist der Verlust genau benennbar -- es ist die
Windungszahl, und zwar eine ganze Zahl je Richtung. Das ist der einzige Ort der
Serie, an dem ein Informationsverlust \emph{exakt} quantifiziert werden kann,
und er ist es, weil die verlorene Größe ganzzahlig ist.

\section{Prüfskript}

\begin{itemize}
  \item Ganzzahligkeit, Koordinatenfreiheit und Skalenunabhängigkeit von
    $\Delta w$; Gegenprobe, dass eine zellbasierte Entropie von der
    Einteilung abhängt und eine spektrale nicht:
    \texttt{python/A\_Serie\_Skripte/a180\_windungsquant.py}
\end{itemize}


\section{Zeichenerklärung}

Symbole die in der gesamten A-Serie verwendet werden, sind in A010 erklärt.
Spezifisch für dieses Dokument:

\begin{center}
\renewcommand{\arraystretch}{1.3}
{\small\begin{tabular}{lp{9cm}}
\toprule
Symbol & Bedeutung \\
\midrule
$S_L = k_B\ln 2$ & Landauer-Energie; Mindestaufwand pro gelöschtem Bit \\
$\Delta w$ & Wicklungsänderung; topologische Diskretheit der Informationsänderung \\
$I_\text{max}$ & Maximale Information bei gegebener Energie \\
$E_L = k_BT\ln 2$ & Landauer-Schranke bei Temperatur $T$ \\
\bottomrule
\end{tabular}}
\renewcommand{\arraystretch}{1.0}
\end{center}

\section{Was offen bleibt}

\textbf{Erstens setzt die Ableitung voraus, dass Bit-Übergang und Windungsquant
identifiziert werden dürfen.} Das ist die physikalische Kernbehauptung; sie wird
nicht weiter reduziert. Die Ableitung ist in A180 ausgeführt (oben), aber die
Behauptung, dass genau $\Delta w=1$ einen Bit trägt, ist eine Setzung -- nicht
aus der Topologie allein zwingend.

\textbf{Zweitens ist nicht gezeigt, dass $\Delta w=1$ die \emph{einzige}
invariante Einheit ist.} Gezeigt ist, dass sie invariant ist und dass
zellbasierte Alternativen es nicht sind. Das ist weniger.

\textbf{Drittens fehlt eine Entropieformel.} Es gibt in dieser Serie keine
ausgeschriebene Entropie als Funktion des Spektrums -- nur die Feststellung, dass
eine solche automatisch invariant wäre. Das ist ein Programm, kein Ergebnis.

\section{Ersetzt}

Dieses Dokument nimmt auf: Dok.~243, 255, 257 (Information), Dok.~290
(Betrag/Phase), Dok.~301 (Windungsquant), Dok.~302 (Elementarzellen und
Partitionsinvarianz).

\section{Verwendung von KI-Werkzeugen}

Dieses Dokument ist unter Verwendung KI-gestützter Werkzeuge entstanden. Die
Fragestellungen, die Setzungen und alle inhaltlichen Entscheidungen stammen vom
Autor; die Werkzeuge formulieren aus, rechnen nach und erstellen Prüfskripte.
Die Verantwortung für Inhalt und Fehler liegt vollständig beim Autor. Der
ausführliche Wortlaut steht in A010.
